\begin{center}
\SourceLeaf{031}\pdfbookmark[2]{13.}{article-013}13.
\end{center}
\emph{Exemplum.} Sit excentricitas eadem quae in exemplo art. 10,
\[
M = 332^0\,28^{\prime}\,54^{\prime\prime}{,}76.
\]
Hic igitur est $\log e$ in secundis $4{,}704\,1513$, adeoque $e=50600^{\prime\prime}=14^0\,3^{\prime}\,20^{\prime\prime}$. Quare quum hic $E$ minor esse debeat quam $M$, statuemus ad calculum primum $\epsilon=326^0$, unde per tabulas minores fit
\begin{center}
\small
\begin{tabular}{@{}llp{13em}@{}}
\makebox[7em][l]{$\log\sin\epsilon$\hfill .\,.\,.\,.\,.} & $9{,}74756_n,$ & mutatio pro $1^{\prime}$ .\,.\,.\,.\,. $19$, unde $\lambda=0{,}32$\\[.3em]
\makebox[7em][l]{$\log e$ in sec\hfill .\,.\,.} & $4{,}70415$ & \\
\cline{1-2}
 & $4{,}45171_n$ &
\end{tabular}
\end{center}
\begin{center}
\small
\begin{tabular}{@{}llp{12em}@{}}
\multicolumn{2}{@{}l}{hinc $e\sin\epsilon=-28295^{\prime\prime}=-7^0\,51^{\prime}\,35^{\prime\prime}$.} & Mutatio logarithmi pro unitate tabulae, quae hic $10$\\
$M+e\sin\epsilon=$ & $324\ \ 37\ \ 20$ & secundis aequivalet, .\,.\,.\,.\,. $16$; unde $\mu=1{,}6$\\[.3em]
Differt ab $\epsilon$ .\,.\,.\,.\,.\,.\,.\,.\,.\,.\,.\,. & $1\ \ 22\ \ 40=4960^{\prime\prime}$. &
\end{tabular}
\end{center}
\noindent Hinc $\frac{0{,}32}{1{,}28}\times4960^{\prime\prime}=1240^{\prime\prime}=20^{\prime}\,40^{\prime\prime}$. Quare valor correctus ipsius $E$ fit $=324^0\,37^{\prime}\,20^{\prime\prime}-20^{\prime}\,40^{\prime\prime}=324^0\,16^{\prime}\,40^{\prime\prime}$, cum quo calculum secundum tabulas maiores repetemus.
\begin{center}
\begin{tabular}{@{}lll@{}}
\makebox[9em][l]{$\log\sin\epsilon$\hfill .\,.\,.\,.\,.\,.} & $9{,}766\,3058_n$ & $\lambda=29{,}25$\\[.3em]
\makebox[9em][l]{$\log e$\hfill .\,.\,.\,.\,.\,.\,.\,.} & $4{,}704\,1513$ & \\
\cline{1-2}
 & $4{,}470\,4571_n$ & $\mu=147$
\end{tabular}
\end{center}
\begin{center}
\begin{tabular}{@{}ll@{}}
\multicolumn{2}{l}{$e\sin\epsilon=-29543^{\prime\prime}{,}17=-8^0\,12^{\prime}\,23^{\prime\prime}{,}17$}\\[.3em]
$M+e\sin\epsilon=$ & $324\ \ 16\ \ 31{,}59$\\[.3em]
Differt ab $\epsilon$ .\,.\,.\,.\,.\,.\,.\,.\,.\,.\,.\,.\,.\,.\,.\,. & $8{,}41$.
\end{tabular}
\end{center}
\noindent Multiplicata hac differentia per $\frac{\lambda}{\mu-\lambda}=\frac{29{,}25}{117{,}75}$, prodit $2^{\prime\prime}{,}09$, unde valor denuo correctus ipsius $E=324^0\,16^{\prime}\,31^{\prime\prime}{,}59-2^{\prime\prime}{,}09=324^0\,16^{\prime}\,29^{\prime\prime}{,}50$, intra $0^{\prime\prime}{,}01$ exactus.
